\section{Why This Folder Exists}

Imagine you have already solved the ``easy'' version of this problem: for a
single pinhole, what diameter gives the sharpest possible image for someone
with a given amount of short-sightedness? That question was answered earlier
in this project, and the answer changes with myopia power $M$ (in diopters)
--- a more short-sighted eye needs a smaller hole. That single-hole answer is
stored in \texttt{new\_sweeps/per\_M\_optimum\_map.csv} and is treated as a
\emph{fixed input} to everything in this handbook: we are not re-deriving hole
size here.

The harder question is: sunglasses need to block most of the light, and a
single tiny hole lets almost none through. The natural fix is to use
\emph{many} small holes instead of one. But the moment you have more than one
hole, a brand-new problem appears that a single hole never has: the eye can
form more than one image of the same scene at once, one through each hole.
If the eye is not perfectly focused (which, for someone wearing tinted
glasses rather than corrective lenses, is basically always true to some
degree), those images land in \emph{slightly different places} on the
retina. A small offset just looks like extra blur. A large offset looks like
literal double vision --- a real, disqualifying defect.

This handbook is about the code that answers one specific question:
\textbf{given a fixed hole size (already chosen to be optimal), how should
those many holes be arranged so that the double-vision problem never
happens, for every level of myopia and every direction the eye might be
looking?}

\subsection{The Two Physical Effects That Matter}

\intuitbox{Two physically independent things can go wrong with a multi-hole array, and
the code treats them as two separate checks, not one:
\begin{enumerate}
  \item \textbf{Ghosting.} If two holes are far enough apart, and the eye is
  out of focus, the two copies of the image they produce land in different
  places on the retina --- literally two overlapping pictures.
  \item \textbf{Vignetting.} As the eye rotates to look around a scene, the
  pupil --- a fixed-size circular opening --- physically slides across the
  pattern of holes. Holes that were far from the pupil edge at straight-ahead
  gaze can suddenly come into play once the eye turns, changing \emph{which}
  holes are simultaneously active and therefore changing the ghosting
  picture entirely.
\end{enumerate}
A pattern that is perfectly safe when the eye looks straight ahead can become
unsafe the moment the eye looks 15 degrees to one side. This is the single
most important finding in the whole study, and it is why the code always
checks both effects separately rather than assuming one implies the other.
}
\subsection{What the Code Does, in One Sentence}

For a chosen arrangement of holes (a \emph{geometry}), at a chosen hole
count and spacing, the code computes: (a) will any two holes ever produce a
visible double image, at any myopia level, and (b) will that stay true as
the eye looks around --- and if both answers are ``no'', the design is
\emph{structurally safe}. It then compares many such safe designs against
each other to find the one that packs the most holes (and therefore lets the
most light through) without crossing either safety line.

\newpage
